\documentstyle[%


righttag%


,11pt%


,amssymb%


,russian%               remove in Eng.


,izvru%                 Set izven% in Eng.


]{amsart}





%





\newcommand{\tts}[1]{}





%\hoffset=-15mm%                  For Eng.


\setcounter{page}{29}


\god{2003}


\nomer{1~(488)} %                        Remove in Eng.


%\nomer{Vol.\,47, No.\,1}%               Set in Eng.


%\str{\pageref{begin}--\pageref{end}}%  Set in Eng.


\udk{519.63 }





\author{\it П.Н.\,Вабищевич}


% NB:   ^^ remove in Eng.


\title{Численное решение задачи


идентификации правой части параболического уравнения}


\thanks{Работа выполнена при финансовой поддержке Российского фонда 


фундаментальных исследований


(проекты \No \No~01-01-00477, 02-01-00555).}





\date{}


%\pagestyle{myheadings}%         Set in Eng.


\pagestyle{plain} %              Remove in Eng.


\begin{document}


%\thispagestyle{plain}%          Set in Eng.


\maketitle


\label{begin}











В теории обратных задач теплообмена \cite{b1}--\cite{b3} рассматриваются


задачи по восстановлению неизвестных источников тепла по


дополнительным измерениям температуры.


Аналогично формулируются и некоторые важные прикладные проблемы


гидрогеологии. 





Можно выделить как самостоятельные задачи случаи, когда решение известно в 


небольшом числе точек \cite{d1}--\cite{d3}. 


В этих условиях для единственности решения необходимо сузить класс


допустимых правых частей параболических уравнений. Во многих


прикладных проблемах естественно считать, что неизвестной является 


правая часть как функция от времени.





Отмеченные обратные задачи в ряде важных случаев удобно


формулировать как задачи управления


\cite{b4}. Для приближенного восстановления


неизвестной правой части используются различные подходы, которые


базируются прежде всего на методах регуляризации \cite{b5}. Такой


общий вычислительный алгоритм при решении существенно некорректных


задач идентификации для многомерных параболических


уравнений применяется, например, в \cite{b6}, \cite{g1}.








В данной работе рассматривается обратная задача восстановления правой части


одномерного параболического уравнения по известному решению.


Отмечаются особенности такой задачи, связанные с ее эволюционностью,


с возможностью последовательного определения правой части


при возрастании времени. Численное решение проводится с


использованием градиентных итерационных методов минимизации


функционала невязки \cite{d4}, \cite{d5}. Регуляризация в итерационных методах


решения некорректных задач достигается за счет согласования числа


итераций с погрешностью входных данных \cite{d5}--\cite{d7}. Представлены


результаты численного решения модельной задачи, которые


демонстрируют возможности идентификации правой части параболического


уравнения при


различном уровне погрешностей во входных данных. 


\medskip





{\bf 1.} {\it Модельная задача}.


В качестве модельной рассмотрим задачу восстановления правой части


одномерного параболического уравнения. Начнем с постановки прямой


задачи. Решения $u(x,t)$ определяются в прямоугольнике


$$


\overline{Q}_T=\overline{\Omega}\times[0,T],


\quad \overline{\Omega}=\{x \mid 0\le x\le l\},


\quad 0\le t\le T.


$$


Функция $u(x,t)$ удовлетворяет уравнению


\begin{equation}\label{6.45}


\frac{\partial u}{\partial t}=


\frac\partial{\partial x}\bigg(k(x)\frac{\partial u}{\partial


x}\bigg) + f(x,t) , \quad 0 < x < l, \quad 0 < t \le T ,


\end{equation}


при стандартных ограничениях $k(x) \geq \varkappa >0$.


Граничные и начальные условия для простоты возьмем однородными


\begin{gather}\label{6.46}


u(0,t)=0, \quad u(l,t)=0, \quad 0 < t \le T , \\


%


\label{6.47}


u(x,0)=0, \quad 0 \le x \le l .


\end{gather}





В прямой задаче (\ref{6.45})--(\ref{6.47}) решение $u(x,t)$


определяется по известным коэффициенту $k(x)$ и правой части


$f(x,t)$. В рассматриваемой обратной задаче неизвестной будет


правая часть $f(x,t)$ (мощность источников), но само решение


$u(x,t)$ предполагается известным. Для вычисления правой части


имеем из (\ref{6.45}) явную формулу


\begin{equation}\label{6.48}


f(x,t) = \frac{\partial u}{\partial t} -


\frac\partial{\partial x}\bigg(k(x)\frac{\partial u}{\partial


x}\bigg) , \quad 0 < x < l, \quad 0 < t \le T .


\end{equation}





Входные данные заданы с погрешностью, и поэтому прямое


использование (\ref{6.48}) затруднительно.  Наиболее


принципиальным является влияние неточного задания решения. Пусть


вместо точного решения задачи (\ref{6.45})--(\ref{6.47}) $u(x,t)$


известна функция $u_{\delta}(x,t)$, причем в некоторой норме


параметр $\delta$ определяет уровень погрешности в задании


решения, т.\,е.


\begin{equation}\label{6.49}


\|u_{\delta}(x,t) - u(x,t)\| \le \delta .


\end{equation}


Нужно использовать специальные вычислительные алгоритмы


устойчивого численного дифференцирования. Для приближенного


решения можно использовать алгоритмы конечноразностной


регуляризации, когда в качестве параметра регуляризации выступают


шаги дискретизации. Мы не будем  подробно обсуждать возможности такого подхода


в данной работе. Отметим только, что в этом случае при решении обратной


задачи (\ref{6.46})--(\ref{6.49}) необходимо согласовывать с


погрешностью в задании  $u(x,t)$ шаг дискретизации по пространству


и по времени.





Эволюционные  задачи обладают определенной и весьма существенной


спецификой. Решение задачи на текущий момент времени зависит


только от предыстории, от решения на предшествующие моменты


времени и не зависит от решения на последующие моменты времени.


Такая особенность эволюционных задач может (и часто просто


должна) учитываться при построении вычислительных алгоритмов. Это


абсолютно оправдано при рассмотрении прямых задач типа


(\ref{6.45})--(\ref{6.47}), это может быть также важно при


рассмотрении обратных задач, подобных (\ref{6.46})--(\ref{6.49}).





При решении эволюционных задач в  общем контексте можно


говорить о двух типах вычислительных алгоритмов. Первые из них


основаны на определении решения по решению на предшествующие


моменты времени. В этом случае мы будем говорить о


{\it локальных алгоритмах решения эволюционных


задач}. В


{\it глобальных алгоритмах} для нахождения решения на заданный


момент задействованы будущие моменты времени.


\medskip





{\bf 2.} {\it Глобальная регуляризация}.


Будем использовать общую схему регуляризации по А.Н.\,Ти\-хо\-но\-ву для


приближенного решения обратной задачи (\ref{6.46})--(\ref{6.49}).





В гильбертовом пространстве $\cal{L}_2(\Omega)$ норму и


скалярное произведение введем обычным образом:


$$


(v,w) = \int_{\Omega} v(x) w(x) d x ,


\quad


\|v\|^2 = (v,v) = \int_{\Omega} v^2(x) d x .


$$


Для функций $v(x,t), w(x,t) \in \cal{H}$, где $\cal{H} =


\cal{L}_2(Q_{T})$, положим


$$


(v,w)_* = \int_0^T (v,w) dt = \int_0^T \int_{\Omega} v(x) w(x) d x  dt,


\quad \|v\|_{*} = (v,v)_*^{1/2} .


$$


Определим оператор


$$


\cal{A} u =  - \frac{d}{dx} \bigg ( k(x) \frac{du}{dx} \bigg )


$$


на множестве функций, удовлетворяющих (\ref{6.46}). В


$\cal{L}_2(\Omega)$


$$


\cal{A} = \cal{A}^* \ge m E,


\quad m = \varkappa \frac{\pi^2}{l^2} .


$$





Обратная задача (\ref{6.46})--(\ref{6.48}) записывается в виде


\begin{equation}\label{6.50}


f =\cal{D} u ,


\end{equation}


где оператор


\begin{equation}\label{6.51}


\cal{D} u = \frac{d u}{d t} + \cal{A} u


\end{equation}


определен на множестве функций, удовлетворяющих начальному условию


(\ref{6.47}). Для входных данных (см. (\ref{6.49})) имеем


\begin{equation}\label{6.52}


\|u_{\delta} - u \|_* \le \delta .


\end{equation}








При использовании метода регуляризации А.Н.\,Тихонова приближенное


решении $f_{\alpha}$ задачи (\ref{6.50})--(\ref{6.52})


определяется как решение вариационной задачи


\begin{equation}\label{6.53}


(E + \alpha \cal{D}^{*} \cal{D}) f_{\alpha} =


\cal{D}u_{\delta} .


\end{equation}


Специфика применения метода регуляризации к решению


рассматриваемой обратной эволюционной задачи проявляется в


операторе задачи $\cal{D}$, определяемом согласно


(\ref{6.51}). В частности, необходимо явно определить оператор


$\cal{D}^{*}$.





Имеем


$$


(\cal{D} v, w)_* = \int_0^T \bigg( \frac{d v}{d t},


w \bigg) dt + \int_0^T (\cal{A}, w) dt =


(v, w) \big|_0^T - \int_0^T \bigg( v, \frac{d w}{d t},


w \bigg) dt + \int_0^T (\cal{A}, w) dt = (v, \cal{D}


w)_*


$$


при условии, что $v(x,0)$, а $w(x,T) = 0$. Тем самым оператор


$\cal{D}^{*}$ определяем как


\begin{equation}\label{6.54}


\cal{D}^* w = - \frac{d w}{d t} + \cal{A}w


\end{equation}


на множестве функций


\begin{equation}\label{6.55}


w(x,T)=0, \quad 0 \le x \le l .


\end{equation}








С учетом (\ref{6.47}), (\ref{6.51}), (\ref{6.54}), (\ref{6.55})


уравнение (\ref{6.53}) для нахождения приближенного решения


$f_{\alpha}$ есть эллиптическое уравнение, которое включает вторые


производные по времени и четвертые по пространству. Для


$f_{\alpha}$ реализуются следующие граничные условия при $t=0,T$:


\begin{alignat}{2}\label{6.56}


f_{\alpha} (x,0) &= 0, &\quad 0 &\le x \le l, \\


%


\label{6.57}


\cal{D} f_{\alpha} (x,T) &= 0, &\quad 0 &\le x \le l .


\end{alignat}


Эти моменты должны аккуратно отслеживаться при численной


реализации рассматриваемого подхода.





При рассмотрении задачи идентификации правой части 


дифференциального уравнения второго порядка можно отметить две основные


возможности. Первая из них связана с применением  схемы


регуляризации А.Н.\,Тихонова при интерпретации 


задачи идентификации правой части как задачи решения операторного


уравнения первого рода, вторая


--- как задачи вычисления значений неограниченного оператора.


Эта последняя возможность реализована в схеме (\ref{6.52}),


(\ref{6.53}) при приближенном решении задачи (\ref{6.46})--(\ref{6.49}).


Имеет смысл также остановиться на стандартном


варианте регуляризации А.Н.\,Тихонова.





По заданной правой части $f(x,t)$ из решения краевой задачи


(\ref{6.45})--(\ref{6.47}) однозначно определяется решение


$u(x,t)$. Это соответствие отобразим введением оператора


\begin{equation}\label{6.58}


\cal{G} f = u.


\end{equation}


Вместо $u(x,t)$ задана функция $u_{\delta}(x,t)$, причем выполнена


оценка (\ref{6.52}).





Приближенное решение $f_{\alpha}$ задачи (\ref{6.49}),


(\ref{6.58}) определяется как решение задачи


\begin{equation}\label{6.59}


J_{\alpha}(f_{\alpha}) = \min_{v \in \cal{H}} J_{\alpha} (v),


\end{equation}


где


\begin{equation}\label{6.60}


J_{\alpha} (v) = \|\cal{G} v - u_{\delta} \|_*^2 + \alpha


\|v\|_*^2 .


\end{equation}





Необходимо обратить внимание на следующее важное обстоятельство. При


использовании алгоритма (\ref{6.52}), (\ref{6.53}), построенного


на основе  интерпретации задачи идентификации как задачи


вычисления значений неограниченного оператора, ставятся


дополнительные ограничения на искомую функцию (граничные условия


типа (\ref{6.56}), (\ref{6.57})). 


Это часто не совсем оправдано при приближенном


решении прикладных проблем.


При использовании (\ref{6.59}), (\ref{6.60}) такие проблемы не


возникают.


\medskip





{\bf 3.} {\it Локальная регуляризация}.


При приближенном решении задачи идентификации правой части


нестационарного уравнения по известному решению часто удобнее


ориентироваться на алгоритмы  определения правой части на заданный


момент времени с использованием входной информации только на


предшествующие моменты времени. По сравнению с алгоритмами


глобальной регуляризации в этом случае, вообще говоря, теряем в


точности приближенного решения, но приобретаем в плане


оперативности \cite{g2}. Отметим некоторые основные возможности по


построению алгоритмов локальной регуляризации для приближенного


решения обратной задачи (\ref{6.46})--(\ref{6.49}). 


Oстановимся на локальном аналоге регуляризации типа (\ref{6.59}),


(\ref{6.60}).





Основная идея связана с тем, что правая часть определяется по


решению на каждый фиксированный момент. Другими словами,


регуляризация процедуры численного дифференцирования проводится


только по пространственным переменным. Фактически сглаживаются


входные данные только по части переменных. Такого типа алгоритмы


реализуются при предварительной дискретизации по времени.





Введем равномерную сетку по времени


$$


\ov\omega_\tau =


\omega_\tau\cup \{T\} =


\{t_n=n\tau,


\ \; n=0,1,\dotsc,N_0,


\ \; \tau N_0=T\} .


$$


Будем пользоваться следующими основными безындексными обозначениями


теории разностных схем \cite{b12}:


$$


y=y^n, \ \ \wh y=y^{n+1}, \ \ \check y=y^{n-1},


\ \ y_{\ov t}=\frac {y-\check y}{\tau},


\ \ y_t=\frac {\wh y-y}{\tau}.


$$





Сформулируем обратную задачу идентификации правой части


параболического уравнения после частичной дискретизации. Для


основных величин сохраним те же обозначения, что и для


непрерывного случая. Для простоты ограничимся чисто


неявной аппроксимацией по времени, когда правая часть определяется


(ср. с (\ref{6.48})) из дифференциально-разностного соотношения


\begin{equation}\label{6.61}


f^n =\frac{u^n - u^{n-1}}{\tau} + \cal{A} u^n,


\quad n = 1,2,\dotsc, N_0.


\end{equation}


Входные данные (решение прямой задачи $u^n$) заданы с


погрешностями. Будем считать, что уровень погрешностей определяется


постоянной $\delta$, причем


\begin{equation}\label{6.62}


\|u^n_{\delta} - u^n\| \le \delta ,


\quad n = 1,2,\dotsc, N_0.


\end{equation}





Приближенное решение задачи (\ref{6.61}), (\ref{6.62}) на момент


времени $t = t_n$ обозначим $f_{\alpha}^n$. Рассмотрим вопрос о


восстановлении функции $f_{\alpha}^n$ по заданным $u^n_{\delta}$,


$u^{n-1}_{\delta}$, $f_{\alpha}^{n-1}$. Будем считать, что


приближенной правой части  $f_{\alpha}^n$ соответствует решение


краевой задачи


\begin{gather}\label{6.63}


\frac{w^n - w^{n-1}}{\tau} + \cal{A} w^n = f_{\alpha}^n,


\quad n = 1,2,\dotsc, N_0, \\


%


\label{6.64}


w^0 = 0 .


\end{gather}





Сначала определим $f_{\alpha}^n$ как решение задачи


%\begin{equation}\label{6.65}


$$


J_{\alpha}^n (f_{\alpha}^n) = \min_{v \in \cal{H}}


J_{\alpha}^n (v),


%%\end{equation}


$$


где $\cal{H} = \cal{L}_2(\Omega)$   и


\begin{equation}\label{6.66}


J_{\alpha}^n (v) =


\|u_{\delta} - w(v) \|^2 + \alpha \|v\|^2 .


\end{equation}


Здесь с учетом (\ref{6.63}), (\ref{6.64}) $w(v)$ есть решение


сеточной задачи


\begin{equation}\label{6.67}


\frac{w - \check{w}}{\tau} + \cal{A} w = v .


\end{equation}


Это позволяет записать (\ref{6.66}) в виде


$$


J_{\alpha}^n (v) =


\bigg \|u_{\delta} - \frac{1}{\tau} \cal{G}_{\tau}\check{w} -


\cal{G}_{\tau}v \bigg \|^2 + \alpha \|v\|^2 ,


$$


где


$$


\cal{G}_{\tau} = \cal{G}_{\tau}^* = \bigg (\frac{1}{\tau} E +


\cal{A} \bigg)^{-1} .


$$


Тем самым приближенное решение на каждый момент времени $t = t_n$


определяется из уравнения


%%\begin{equation}\label{6.68}


$\cal{G}_{\tau}^* \cal{G}_{\tau} f_{\alpha} +


\alpha f_{\alpha} = \cal{G}_{\tau}^* u_{\delta} - \frac{1}{\tau}


\cal{G}_{\tau}^* \cal{G}_{\tau}\check{w}$.


%%\end{equation}


С учетом (\ref{6.67}) и введенных обозначений получим задачу


%%\begin{equation}\label{6.69}


$\cal{G}_{\tau}^{-1} w = f_{\alpha} + \frac{1}{\tau}


\check{w}$


%%\end{equation}


для определения $w_n$.


\medskip





{\bf 4.} {\it Итерационнoe решение задачи идентификации}.


При решении обратных задач математической физики наибольшего


внимания заслуживают итерационные методы, которые наиболее четко


реализуют идею нахождения решения обратной задачи через


последовательное решение набора прямых задач. При приближенном


решении задачи (\ref{6.46})--(\ref{6.49}) итерационными методами


остановимся на случае глобальной регуляризации в варианте


(\ref{6.58})--(\ref{6.60}).





После симметризации уравнения (\ref{6.58}) двухслойный


итерационный метод запишем в виде


\begin{equation}\label{6.70}


\frac{f_{k+1} - f_k}{\tau_{k+1}} + \cal{G}^* \cal{G} f_k


= \cal{G}^* u_{\delta} ,


\quad  k = 0,1, \dotsc


\end{equation}


Итерационные параметры при использовании метода скорейшего спуска


рассчитываются по формуле


\begin{equation}\label{6.71}


\tau_{k+1} =  \frac{\|r_k\|_*^2}{\|\cal{G} r_k\|_*^2} ,


\quad r_k =\cal{G}^* \cal{G} f_k - \cal{G}^* u_{\delta} ,


\quad k =0,1, \dotsc


\end{equation}


Число итераций в (\ref{6.70}), (\ref{6.71}) согласуется с


погрешностью $\delta$ (см. (\ref{6.52})).





Реализация такого подхода связана с возможностью вычисления


значений оператора $\cal{G}$ и $\cal{G}^*$. Напомним, что


$v = \cal{G} f_k$ есть решение прямой задачи


\begin{gather}\label{6.72}


\frac{d v}{d t} + \cal{A} v = f_k ,


\quad 0 < t \le T, \\


%


\label{6.73}


v(0) = 0 .


\end{gather}


При нахождении значений сопряженного оператора $w = \cal{G}^*


v$ решается прямая задача


\begin{gather}\label{6.74}


- \frac{d w}{d t} + \cal{A} w = v ,


\quad 0 \le t <  T, \\


%


\label{6.75}


 w(T) = 0 .


\end{gather}


Тем самым, при заданном итерационном параметре переход на новую


итерацию в соответствии с (\ref{6.70}) связан с решением двух


прямых задач (\ref{6.72}), (\ref{6.73}) и (\ref{6.74}),


(\ref{6.75}).





Отметим некоторые основные особенности вычислительной реализации


рассматриваемого итерационного метода, которые связаны прежде


всего с дискретизацией по времени. Сохраним за сеточными


функциями те же обозначения, что и для функций непрерывного


аргумента.


Будем использовать на интервале $\ov{\Omega} = [0,l]$ с шагом $h$


равномерную сетку \tts{on the interval with the step? \%-( }


$$


\ov{\omega} = \{x\mid x = x_i  = ih,


\ i = 0,1,\dotsc,N,\ Nh = l\},


$$


где $\omega$ --- множество внутренних узлов, а $\partial\omega$


--- множество граничных узлов.





Во внутренних узлах дифференциальный оператор $\cal{A}$


аппроксимируем разностным оператором второго порядка \cite{b12}


$$


%%\begin{equation}\label{6.76}


A y = - (ay_{\ov{x}})_x , \quad x \in \omega ,


%% \end{equation}


$$


где, например, $a(x) = k(x-0.5h)$.





В сеточном гильбертовом пространстве $L_2(\omega)$ норму введем


соотношением $\|y\| = (y,y)^{1/2}$, где


$(y,w) = \sum\limits_{x\in \omega} y(x) w(x) h$.


Напомним, что на множестве функций, обращающихся в нуль на


$\partial\omega$, для самосопряженного  оператора $A$ при


ограничениях $k(x) \ge \varkappa


> 0$, $q(x) \ge 0$ верна оценка


$$


%% \begin{equation}\label{6.77}


A = A^* \ge \varkappa \lambda_0 E,


%% \end{equation}


$$


где


$$


\lambda_0 = \frac {4} {h^2}


\sin^2 \frac {\pi h}{2l} \ge \frac {8}{l^2} .


$$





Прямой задаче (\ref{6.72}), (\ref{6.73}) поставим в соответствие


симметричную разностную задачу


\begin{gather}\label{6.78}


\frac{v^{n} - v^{n-1}}{\tau} + \frac{1}{2}A (v^{n+1} + v^{n}) =


\frac{1}{2} (f_k^{n+1} + f_k^{n}) ,


\quad n = 1,2,\dotsc, N_0, \\


%


\label{6.79}


v^0 = 0 ,


\quad x \in \omega .


\end{gather}


В этом случае погрешность аппроксимации имеет второй порядок как


по времени, так и по пространству. Аналогично рассматриваются и


другие двухслойные разностные схемы. Задаче (\ref{6.78}),


(\ref{6.79}) в операторной записи придадим форму $v = G f_k$,


которая определяет оператор $G$.





Для двумерных сеточных функций определим гильбертово пространство


$H = L_2(Q_{T})$, в котором скалярное произведение и норма


определены следующим образом:


$$


(v,w)^* = \sum_{n=1}^{N_0 -1} (v^n, w^n) \tau +


\frac{\tau}{2} (v^0, w^0) +


\frac{\tau}{2} (v^{N_0}, w^{N_0}) ,


\quad \|v\|^* = \sqrt{(v,v)^*} .


$$





Сопряженной к (\ref{6.78}), (\ref{6.79}) в $H$ будет сеточная


задача (см. (\ref{6.74}), (\ref{6.75}))


\begin{gather}\label{6.80}


- \frac{w^{n} - w^{n-1}}{\tau} + \frac{1}{2}A(w^{n} + w^{n-1}) =


\frac{1}{2}( v^{n} + v^{n-1}),


\quad n = 1,2,\dotsc, N_0, \\


%


\label{6.81}


w^{N_0} = 0 ,


\quad x \in \omega .


\end{gather}


В этом можно убедиться, если домножить скалярно уравнение


(\ref{6.78}) на $w^n$. Задаче (\ref{6.80}), (\ref{6.81})


соответствует компактная запись $w = G^* v$.





В соответствии с (\ref{6.70}) итерационный метод записывается в


виде


\begin{equation}\label{6.82}


\frac{f_{k+1} - f_k}{\tau_{k+1}} + G^* G f_k


= G^* u_{\delta} ,


\quad  k = 0,1, \dotsc


\end{equation}


На первом этапе проводится расчет правой части $G^* u_{\delta}$


--- с этой целью решается краевая задача типа (\ref{6.80}),


(\ref{6.81}). Вычисление невязки $r_k = G^* G f_k - G^*


u_{\delta}$ требует на каждом шаге решения двух краевых задач


((\ref{6.78}), (\ref{6.79}) и (\ref{6.80}), (\ref{6.81})). Расчет


итерационных параметров проводится (см. (\ref{6.71})) по формуле


$$


%%\begin{equation}\label{6.83}


\tau_{k+1} =  \frac{\|r_k\|_*^2}{\|G r_k\|_*^2} ,


\quad k =0,1, \dotsc


%%\end{equation}


$$


Определение $G r_k$ связано с решением дополнительной краевой


задачи типа (\ref{6.78}), (\ref{6.79}).





В качестве параметра регуляризации выступает число итераций


$K(\delta)$ в (\ref{6.82}). Критерием выхода из итерационного


процесса является


%%\begin{equation}\label{6.84}


$\|G f_{K(\delta)} - u_{\delta}\|_* \le \delta$.


%% \end{equation}





Отметим, что рассматриваемый алгоритм используется при


определенных ограничениях на правую часть. В соответствии с


применяемой симметризацией за счет оператора $G$ на правую часть


накладываются ограничения


$f_{\alpha}^{N_0} = 0$, $x \in \omega$,


$f_{\alpha}^{n} = 0$, $n = 0,1, \dotsc, N$, $x \in \partial\omega$.


\medskip





{\bf 5.} {\it Результаты расчетов}.


Расчеты выполнялись на равномерной сетке с числом узлов $N = 100$,


$N_0 = 100$, когда расчетная область представляла собой единичный


квадрат ($l = 1$, $T = 1$). Обратная задача решалась в рамках


квазиреального эксперимента при


$$


k(x) = 1, \quad f(x,t) = 10 t(1-t) x(1-x) .


$$





Решение задачи при уровне погрешностей $\delta = 0.001$  показано


на рис.~1 (число итераций равно~3), где изображены линии


уровня с шагом $\Delta = 0.1$ точного (штриховые линии) и


приближенного решений. Влияние погрешности на точность


восстановления правой части иллюстрируется на


рис.~2, 3, где представлены результаты решения


с большими и меньшими погрешностями во входных данных. Для решения


задачи с $\delta = 0.0001$ потребовалось сделать 23 итерации.





\begin{figure}[h]


\begin{center}


\begin{picture}(470,160)(0,-15)


\put(120,150){\special{em:graph gr6-5.bmp}}


%% \put(212,-10){\mbox Рис.\,6.}


\end{picture}


\end{center}


%\caption{Рис.\,1. Решение задачи идентификации при  $\delta = 0.001$}


\label{f-6.5}


\end{figure}


\begin{center}


{Рис.\,1. Решение задачи идентификации при  $\delta = 0.001$}


\end{center}





\begin{figure}[h]


\begin{center}


\begin{picture}(470,160)(0,-15)


\put(120,150){\special{em:graph gr6-6.bmp}}


%% \put(212,-10){\mbox Рис.\,6.}


\end{picture}


\end{center}


% \caption{Решение задачи идентификации при $\delta = 0.0001$}


\label{f-6.6}


\end{figure}


\begin{center}


{Рис.\,2. Решение задачи идентификации при $\delta = 0.0001$}


\end{center}


\newpage





\begin{figure}[h]


\begin{center}


\begin{picture}(470,160)(0,-15)


\put(120,150){\special{em:graph gr6-7.bmp}}


%% \put(212,-10){\mbox Рис.\,6.}


\end{picture}


\end{center}


% \caption{Решение задачи идентификации при $\delta = 0.01$}


\label{f-6.7}


\end{figure}


\begin{center}


{Рис.\,3. Решение задачи идентификации при $\delta = 0.01$}


\end{center}





\begin{figure}[h]


\begin{center}


\begin{picture}(470,160)(0,-15)


\put(120,150){\special{em:graph gr6-8.bmp}}


%% \put(212,-10){\mbox Рис.\,6.}


\end{picture}


\end{center}


% \caption{Решение задачи идентификации при $\delta = 0.001$}


\label{f-6.8}


\end{figure}


\begin{center}


{Рис.\,4. Решение задачи идентификации при $\delta = 0.001$}


\end{center}





Естественно, что точность идентификации существенно зависит от


точного решения. В частности, выше отмечалась необходимость


сужения класса искомых правых частей в связи с выполнением


однородныx условий на частях границы расчетной области. На


рис.~4 показаны результаты решения задачи с правой


частью


$$


f(x,t) = 2 x(1-x)


$$


при $\delta = 0.001$.











\begin{thebibliography}{99}





\bibitem{b1}


Алифанов О.М. {\it Обратные  задачи  теплообмена}. --  М.:


Машиностроение, 1988.





\bibitem{b2}


Бек Дж., Блакуэлл Б., Сент-Клэр Ч.


{\it Некорректные обратные задачи теплопроводности}. -- М.: Мир, 1989.





\bibitem{b3}


Samarskii A.A., Vabishchevich P.N. {\it Computational Heat Transfer}.


-- Chichester: Wiley, 1995.





\bibitem{d1}


Денисов А.М. {\it  Введение в теорию


обратных задач}. -- М.: Изд-во Моск. ун-та, 1994.





\bibitem{d2}


Prilepko A.I., Orlovsky D.G.,


Vasin I.A. {\it Methods for solving inverse problems in mathematical


physics} // Monographs and textbooks in pure and appl. math.


--  Marcel Dekker, 2000. -- V.\,231. 





\bibitem{d3}


Борухов В.T., Вабищевич П.Н.


{\it Численное решение обратной задачи восстановления источника в


параболическом уравнении} //


Матем. моделир. -- 1998. -- Т.\,10. -- \No\,11. -- C.\,93--100.





\bibitem{b4}


Лионс Ж.-Л. {\it Оптимальное управление системами, описываемыми


уравнениями с частными производными}. -- М.: Мир, 1972. -- 416 c.





\bibitem{b5}


Тихонов А.Н., Арсенин В.Я. {\it Методы решения


некорректных задач}. --  М.: Наука, 1986.





\bibitem{b6}


Самарский А.А., Вабищевич П.Н.


{\it Разностные методы решения задач идентификации источника для


параболических задач} // Вестн. Моск. ун-та. Вычисл. матем. и киберн.


-- 1995. -- \No\,1. -- С.\,47--56.





\bibitem{g1}


Vabishchevich P.N. {\it


Solving problems of continuation of unsteady fields based on local


point measurements} //


Large-scale sci. comput. of engineering and


environmental probl. -- Vieweg, 1998. -- P.\,111--117. 





\bibitem{d4}


Васильев. Ф.П. {\it Методы решения


экстремальных задач}. -- М.: Наука, 1981. -- 400 c.





\bibitem{d5}


Алифанов О.М., Артюхин Е.А.,


Румянцев С.В. {\it Экстремальные методы решения некорректных задач}.


-- М.: Наука,  1988.





\bibitem{d6}


Бакушинский А.Б.,  Гончарский А.В.


{\it Итеративные методы решения некорректных задач}. -- М.: Наука, 1989.





\bibitem{d7}


Вайникко Г.М., Веретенников А.Ю.


{\it Итерационные процедуры в некорректных задачах}. -- М.: Наука, 1986.


-- 181 c.





\bibitem{g2}


Осипов Ю.С.,  Васильев Ф.П., Потапов М.М.


{\it Основы метода динамической регуляризации}. -- М.: 


Изд-во Моск. ун-та,  1999.





\bibitem{b12}


Самарский А.А. {\it Теория разностных схем}. --  М.: Наука, 1989. -- 616 c.





\end{thebibliography}





\vskip 2cm





\post{\it Институт математического}{\it Поступила}


\post{\it моделирования Российской}{25.06.2002}


\post{\it Академии наук}{}





\label{end}


\end{document}


